\section{Repeated games}
\label{sec:repeated}

\subsection{One match is a repetition of stage games}

A match is not one game but a sequence of them. Every two ticks each player
chooses one of 12 buttons without seeing the other's choice; that simultaneous
move is the stage game of Section~\ref{sec:design}. The Nash bot solves that
stage game afresh each time it commits: it builds a 12$\times$12 matrix,
finds a mixed equilibrium, samples one action and holds it for $c$ decisions.
It carries no belief about who it is playing and no memory of what happened
before (apart from the score, used only to restart at a goal).

In the language of repeated games, this is a \emph{myopic} strategy. If the
stage game has a unique equilibrium and the repetition is finite, backwards
induction unwinds the whole game: in the last period there is no future to
influence, so the equilibrium is played; knowing that, the same argument
applies to the previous period, and so on \cite{fudenberg1991game}.
A memoryless equilibrium policy therefore already implements the finite-horizon
answer. It also explains the bot's main limitation: it cannot punish a
weak opponent for past play or build a reputation, because the past is not
part of its state. With an infinite (or indefinite) horizon the folk theorem
allows many more outcomes to be sustained by credible threats \cite{aumann1981repeated};
realising any of them needs a strategy that conditions on history.

\subsection{Making history visible}

The runner in \texttt{gtbots/arena.py} exposes the one piece of history this
game hides: the opponent's simultaneous button is not observable, but the
button it pressed \emph{last} decision is. Bots that implement
\texttt{observe(opponent\_action)} are called before every \texttt{act}, with
\texttt{None} on the first decision and an integer $0\ldots11$ afterwards.
This adds no state to the existing bots and keeps the interface unchanged for
those that do not use it. It is enough for a Bayesian bot to update a belief
over opponent types (Hasaan's section) and for history-dependent rules such as
tit-for-tat or trigger strategies \cite{axelrod1984cooperation}: the discount factor is then the probability
that play continues to the next decision, which is $1$ until the match ends.

Commitment is the other repeated-game element. Holding one action for $c$
decisions (6 or 14 ticks in the commitment experiment) changes the repeated
game from ``re-sample a mix every 2 ticks'' to a sequence of short contracts.
It removes the jitter a 50/50 mix would produce, and with \texttt{follow=hold}
it makes the bot's own assumption in the look-ahead true for those $c$
decisions. The cost is slower reaction: a change of opponent behaviour is only
noticed at the next planning point.

\subsection{Experiments}

All numbers below come from \texttt{experiments/tune\_heuristic.py} and the
match runner. Matches are first-to-\emph{target} with \texttt{powerUps=False};
each gets its own seed and the Nash bot is on P1 in half of them and on P2 in
the other half. The runner records the winner, goals, ticks, compute seconds,
the replay and the Nash bot's per-plan log for every match. Win rates are
reported with the binomial 95\% interval where $n$ is the number of matches.

\paragraph{Heuristic tuning.}
The old heuristic scored only the ball, giving 4 distinct rows in the
snapshot payoff matrix. Adding the sign-flipping player terms
(\texttt{pair\_x}, \texttt{ball\_lead}, \texttt{pair\_vx}) raises that to 11
distinct rows and 88 distinct cells on the same snapshot, without breaking
zero-sumness: a mirrored position always gets the exact negative score, and
every score stays inside $\pm 0.5$. The weights were tuned by playing
Nash + heuristic against the scripted Baseline; the table reports mean goal
difference per match (positive is Nash winning).

\begin{center}
\begin{tabular}{l|rrrrr|rr}
\hline
name & $b_x$ & $b_{vx}$ & $p_x$ & $\ell$ & $p_{vx}$ & GD & s \\
\hline
start      & 1.5 & 1.0 &  0.4 & 0.6 &  0.2 & $+4$ & 160.7 \\
ball-only  & 1.5 & 1.0 &  0.0 & 0.0 &  0.0 & $-4$ &  19.8 \\
ball-heavy & 2.0 & 1.5 &  0.0 & 0.0 &  0.0 & $-4$ &  19.3 \\
players    & 1.0 & 0.5 &  1.0 & 1.2 &  0.4 & $0$  &  23.6 \\
lead       & 1.0 & 0.8 &  0.3 & 1.5 &  0.2 & $+4$ &  95.4 \\
defensive  & 1.2 & 0.8 & $-0.5$ & 0.8 & $-0.2$ & $+2$ &  54.7 \\
\hline
\end{tabular}
\end{center}

\noindent Here $b_x$, $b_{vx}$, $p_x$, $\ell$ and $p_{vx}$ are the weights on
\texttt{ball\_x}, \texttt{ball\_vx}, \texttt{pair\_x}, \texttt{ball\_lead} and
\texttt{pair\_vx}; GD is Nash goals minus Baseline goals over the four matches
(all ended at the target, none cut off), and ``s'' is the bot's own compute
time in seconds. The two ball-only variants lose every match, so the player
terms are doing real work: they tell the bot whether it is actually in front
of the ball rather than merely on the same side of the pitch as it. Adding the
player terms changes the snapshot matrix from 4 to 11 distinct rows (88
distinct cells), so the equilibrium is also less degenerate.

\noindent \texttt{start} and \texttt{lead} both go 4--0. The tuning script breaks
the tie by candidate order and keeps \texttt{start}; \texttt{lead} is cheaper
(95.4~s versus 160.7~s) and would be the pick if compute were the second
objective. With only four matches per row these win rates are indicative, not
significant: the Clopper--Pearson 95\% interval for 4/4 is about
$0.40$--$1.00$, which overlaps the interval for 0/4. The reported weights are
therefore a selection among these six candidates, not an optimum over all
weight vectors.

\paragraph{Computation.}
The look-ahead is the dominant cost: one plan is 144$\times$30 ticks of
physics, so a shorter commitment means more plans per match. The runner's
per-match seconds and the Nash bot's log (one row per plan) are saved in
\texttt{results/}, so the cost of each commitment length can be read directly
from the match CSVs rather than estimated.

\subsection{What this says about the project}

The Nash bot is an equilibrium of each stage game, not of the repeated game.
That is a defensible benchmark -- it is safe against any opponent, including
one that watches for weaknesses -- but it leaves money on the table against a
predictable opponent such as the camper or the baseline. The opponent-action
hook and the seeded runner are what let the next bot in the tournament use
history, and the tournament/replicator results then say whether a
history-dependent type survives against the myopic one.
